% Exercises: listings (lstset, lstlisting, \lstinline, \lstinputlisting-free, captions, language styles, escapes).
\documentclass{article}
\usepackage{listings,xcolor}
\lstset{basicstyle=\ttfamily\small,keywordstyle=\color{blue},commentstyle=\color{gray},numbers=left,frame=single,breaklines=true}
\lstdefinelanguage{mini}{morekeywords={let,in,fun},sensitive=true,morecomment=[l]{//}}
\begin{document}
\lstlistingname{Code}\lstlistoflistings
\begin{lstlisting}[language=C,caption={A C program},label=lst:c]
#include <stdio.h>
int main(void) { /* comment */
  for (int i = 0; i < 10; i++) printf("%d\n", i);
  return 0;
}
\end{lstlisting}
\begin{lstlisting}[language=Python,caption=Python]
def f(x):
    return x ** 2  # square
\end{lstlisting}
\begin{lstlisting}[language=mini]
let x = 1 in fun y -> x // add
\end{lstlisting}
Listing~\ref{lst:c} and \lstinline[language=C]|int x = 3;| inline.
\end{document}
